% =====================================================================
%  BAO CAO TIEN DO — CXR-LT 2024
%  DICH:   .\build.ps1                       <- cach de nhat, chay 3 luot + nghiem thu
%          xelatex BAO_CAO_TIEN_DO.tex       <- chay 3 LAN (muc luc + do rong longtable)
%          pdflatex BAO_CAO_TIEN_DO.tex      <- cung duoc, dung T5/vntex
%
%  ⚠ DUNG DUNG `latexmk`: no la script PERL, may nay khong cai Perl nen MiKTeX bao
%    "could not find the script engine 'perl'". Do KHONG phai loi cua file .tex.
%    Ca hai engine tren deu da duoc kiem THAT: 0 loi, 10 trang.
%
%  ⚠ "critical issue: You are running MiKTeX on an unsupported version of Windows"
%    chi la loi than phien (MiKTeX 24.1 khong nhan ra Windows 11 build 26200).
%    PDF van ra dung — cu kiem so trang va so dong "^! " trong .log, dung kiem ma thoat.
%  Nguon so lieu: JOURNAL.md, Event.md, Progress.md, CHECKLIST_STAGE5_OPEN.md,
%                 PLAN_VUOT_0281.md, output cua nb61/64/64b/65/66/67/68.
%  Moi con so trong file nay doc tu log/file, khong go tay tu tri nho.
% =====================================================================
\documentclass[11pt,a4paper]{article}

\usepackage{iftex}
\ifPDFTeX
  \usepackage[utf8]{inputenc}
  \usepackage[T5]{fontenc}
\else
  \usepackage{fontspec}
  \setmainfont{Latin Modern Roman}
  \setmonofont{Latin Modern Mono}
\fi

\usepackage[margin=2.4cm]{geometry}
\usepackage{booktabs}
\usepackage{array}
\usepackage{longtable}
\usepackage{amsmath,amssymb}
\usepackage{xcolor}
% KHONG nap `colortbl`: ban 2026/05/01 v1.0l goi \insert@pcolumn — mot lenh chi co o
% `array.sty` moi. Voi array.sty cu (MiKTeX 24.1) thi CHI CAN nap colortbl la MOI cot
% p{} vo ("Undefined control sequence" + "Missing # in alignment preamble"), ke ca khi
% khong dung \rowcolor. Trang thai duoc the hien bang CHU MAU o cot dau, ben hon ne
% nen mau va con doc duoc khi in den trang.
\usepackage{fancyvrb}
\usepackage{tikz}
\usetikzlibrary{positioning,arrows.meta,calc}
\usepackage{enumitem}
\usepackage[hidelinks]{hyperref}
\usepackage{caption}

\captionsetup{font=small,labelfont=bf}
% Cho TeX noi khoang trang thay vi de chu tran ra le. Khong sua duoc mot TU dai hon
% ca dong — cho do phai tu \allowbreak.
\setlength{\emergencystretch}{3em}
\setlist[itemize]{itemsep=2pt,topsep=3pt,parsep=0pt}
\setlist[enumerate]{itemsep=2pt,topsep=3pt,parsep=0pt}

\newcommand{\DONG}{\textbf{\textcolor{red!62!black}{ĐÓNG}}}
\newcommand{\MO}{\textbf{\textcolor{green!42!black}{MỞ}}}
\newcommand{\XONG}{\textbf{\textcolor{blue!58!black}{XONG}}}
\newcommand{\code}[1]{\texttt{\small #1}}

\title{\textbf{CXR-LT 2024 — Báo cáo tiến độ}\\[2mm]
       \large Kiến trúc, các hướng nghiên cứu, và khoảng cách tới mục tiêu}
\author{Nhóm CXR $\times$ EHR}
\date{10 tháng 9, 2026}

\begin{document}
\maketitle

\begin{center}
\begin{tabular}{@{}lll@{}}
\toprule
\textbf{Chỉ số hiện tại (test)} & \textbf{0.2625} & macro-mAP, 40 nhãn, 78{,}946 ảnh \\
\textbf{Nhà vô địch (đội A)} & 0.281 & CXR-LT 2024 Task 1 \\
\textbf{Khoảng cách} & \textbf{0.0185} & tương đương hạng 7--8 / 17 đội \\
\bottomrule
\end{tabular}
\end{center}

\vspace{2mm}


% =====================================================================
\section{Kiến trúc hiện tại}
% =====================================================================

\subsection{Bài toán và dữ liệu}

Phân loại \textbf{đa nhãn} 40 bệnh lý trên phim X-quang ngực, dữ liệu MIMIC-CXR theo
thiết lập chính thức của thử thách CXR-LT 2024 (Lin et al., arXiv:2506.07984). Thiết
lập của chúng ta khớp tuyệt đối với ban tổ chức, nên con số so sánh trực tiếp được
với bảng xếp hạng.

\begin{table}[h]
\centering
\small
\begin{tabular}{@{}llr@{}}
\toprule
\textbf{Thành phần} & \textbf{Mô tả} & \textbf{Số lượng} \\
\midrule
Train        & chia theo bệnh nhân 70\%          & 258{,}871 ảnh \\
Development  & chia theo bệnh nhân 10\%          & 39{,}293 ảnh \\
Test         & chia theo bệnh nhân 20\%          & 78{,}946 ảnh \\
Nhãn Task 1  & đa nhãn, mất cân bằng nặng        & 40 \\
Nhãn Task 3  & zero-shot, \textbf{không được chạm} & 5 \\
Thước đo     & macro-averaged AP; hoà thì xét mAUROC & --- \\
\bottomrule
\end{tabular}
\caption{Thiết lập dữ liệu chính thức. Tuân thủ nhãn đã được kiểm: 5 nhãn zero-shot
(\code{Bulla}, \code{Cardiomyopathy}, \code{Hilum}, \code{Osteopenia}, \code{Scoliosis})
không xuất hiện trong tập 40 nhãn của ta.}
\end{table}

\paragraph{Phân tầng theo độ hiếm.} Ban tổ chức chia theo \emph{tỉ lệ hiện} trên train:
head $>10\%$, medium $1$--$10\%$, tail $<1\%$, cho \textbf{9 / 15 / 16} nhãn. Thang riêng
mà ta dùng trước ngày 08/09 chia theo \emph{số ca} và cho 13 / 16 / 11 --- mọi báo cáo
lớp trước mốc đó dùng thang khác, cần đối chiếu cẩn thận.

\subsection{Sơ đồ tổng thể}

Hệ thống là \textbf{ensemble hai thành viên} có kiến trúc khác nhau, trộn theo lớp ở
mức xác suất đầu ra.

\begin{center}
\begin{tikzpicture}[
  font=\small,
  node distance=0mm,
  box/.style   = {draw, rounded corners=2pt, align=center, inner sep=4pt,
                  minimum height=8mm, text width=#1},
  bb/.style    = {box=40mm, fill=blue!5},
  head/.style  = {box=33mm, fill=black!4},
  ctx/.style   = {box=30mm, fill=orange!8, dashed},
  io/.style    = {box=50mm, fill=none},
  res/.style   = {box=44mm, fill=green!7},
  ar/.style    = {-{Latex[length=2mm]}, thick},
  lb/.style    = {font=\scriptsize, inner sep=1.5pt},
]

\def\xL{-4.15}   \def\xR{4.15}

% --- hang 1: dau vao ---
\node[io]  (img) at (0, 0) {\textbf{ẢNH X-QUANG NGỰC} \\[-1pt] \scriptsize $518 \times 518$};

% --- hang 2: hai backbone ---
\node[bb]  (A) at (\xL, -2.05) {\textbf{[A] RAD-DINO} \\[-1pt] \scriptsize ViT-B/14 · \textbf{đông cứng}
                                \\[-2pt] \scriptsize 1369 patch · 9.46 mm};
\node[bb]  (G) at (\xR, -2.05) {\textbf{[G] ConvNeXtV2} \\[-1pt] \scriptsize nano · \textbf{có huấn luyện}
                                \\[-2pt] \scriptsize khởi tạo từ PadChest};

% --- hang 3: hai dau doc giong het nhau ---
\node[head] (HL) at (\xL, -4.75) {\textbf{ĐẦU ĐỌC} \\[-1pt]
        \scriptsize [E] lớp1Img \\[-2pt] \scriptsize [D] lớp1
        \\[-2pt] \scriptsize [B] ML-Decoder};
\node[head] (HR) at (\xR, -4.75) {\textbf{ĐẦU ĐỌC} \\[-1pt]
        \scriptsize [E] lớp1Img \\[-2pt] \scriptsize [D] lớp1
        \\[-2pt] \scriptsize [B] ML-Decoder};

% --- ngu canh o giua, cap cho CA HAI ---
\node[ctx] (C) at (0, -4.75) {\textbf{NGỮ CẢNH} \\[-1pt]
        \scriptsize phim chụp trước \\[-2pt] \scriptsize token EHR (\code{ColTokBag})};

% --- hang 4: tron + ket qua ---
\node[box=64mm, fill=black!4] (H) at (0, -7.15) {\textbf{[H] TRỘN THEO LỚP}
        \\[-1pt] \scriptsize head 0.40 \quad medium 0.40 \quad tail 0.00};
\node[res] (R) at (0, -8.75) {\textbf{macro-mAP $=$ 0.2625}};

% --- mui ten ---
\draw[ar] (img.south) -- ++(0,-0.35) -| (A.north);
\draw[ar] (img.south) -- ++(0,-0.35) -| (G.north);
\draw[ar] (A) -- node[lb, right] {patch token} (HL);
\draw[ar] (G) -- node[lb, right] {patch token} (HR);
\draw[ar] (C) -- (HL);
\draw[ar] (C) -- (HR);
\draw[ar] (HL.south) -- ++(0,-0.55) node[lb, above right] {logits 40} -| (H.north);
\draw[ar] (HR.south) -- ++(0,-0.55) node[lb, above left]  {logits 40} -| (H.north);
\draw[ar] (H) -- (R);

\end{tikzpicture}
\end{center}

\noindent Hai nhánh dùng \textbf{cùng một mã đầu đọc} nhưng là hai bản sao huấn luyện độc
lập. Khác biệt thật nằm ở backbone: RAD-DINO \emph{đông cứng} (chỉ 7.3M tham số đầu đọc
được học), còn ConvNeXtV2 thì \emph{huấn luyện cả backbone}. Thứ tự trong đầu đọc là
\code{lớp1Img} (phim trước, không gian \code{D\_CXR}) $\to$ \code{pin} $\to$ \code{lớp1}
(token EHR) $\to$ ML-Decoder.

\subsection{Các module}

\subsubsection*{[A] Nhánh ảnh --- backbone RAD-DINO}

RAD-DINO ViT-B/14 ở độ phân giải $518^2$, cho $37\times37 = 1369$ patch token
($9{,}46$ mm mỗi patch). Backbone được fine-tune một lần ở notebook \code{38a}
(bản \code{a2\_E3}) rồi \textbf{đông cứng hoàn toàn} cho mọi thí nghiệm sau. Việc mở
băng thêm đã được thử và đóng lại (mục~\ref{sec:closed}).

Hệ quả quan trọng của việc freeze: mọi seed của ta \emph{dùng chung} 86M tham số
giống hệt nhau, chỉ khác 7.3M tham số đầu đọc. Đây là lý do việc thêm seed chỉ cho
$+0.0008$, và cũng là lý do hướng \emph{ensemble backbone fine-tune độc lập} vẫn còn mở.

\subsubsection*{[B] ML-Decoder --- đầu đọc đa nhãn}

Mỗi bệnh lý được gán một vector truy vấn (Q) riêng. Bốn mươi vector này giữ nguyên từ lúc
khởi tạo và không tham gia huấn luyện, nên toàn bộ việc học dồn vào phép chú ý và các
lớp phía sau.
Từng truy vấn quét khắp các patch của ảnh rồi gom về một biểu diễn dành riêng cho bệnh
lý đó. Các truy vấn không trao đổi với nhau, nên bốn mươi bệnh lý được đọc song song và
độc lập vì trong bài toán đa nhãn một phim có thể mang nhiều bệnh cùng lúc.
Ở bước chấm điểm, mỗi bệnh lý dùng một vector trọng số của riêng nó trên đúng biểu diễn
của nó, thay vì trải phẳng cả bốn mươi biểu diễn rồi đưa qua một lớp tuyến tính chung.
Số tham số nhờ vậy giảm bốn mươi lần.

\begin{Verbatim}[fontsize=\small,frame=single,framesep=2mm]
q  = class_query[40]                       # cố định, không học
kv = pin(patch)                            # 768 -> 512
kv = lớp1(kv, ein(etok))                  # điều biến bởi EHR   [module D]
h  = XAttn(q, kv)                          # LayerNorm + FFN residual
logits = GroupFC(h)                        # [B, 40]
\end{Verbatim}

\subsubsection*{[C] Bộ mã hoá EHR --- \code{ColTokBag}}

Mỗi lần chụp đi kèm 257 con số: 217 đặc trưng lâm sàng (xét nghiệm, sinh hiệu, bệnh nền,
các cờ vi sinh) và 40 điểm số do một mô hình cây chấm trước đó trên chính bảng EHR.

Thay vì nén cả 257 số thành một vector duy nhất, mỗi cột được cấp một token riêng. Danh
tính token thuộc về \emph{cột}: mỗi đặc trưng có một hướng cố định trong không gian 256
chiều, dùng chung cho mọi bệnh nhân. Giá trị của từng bệnh nhân chỉ quyết định \emph{biên
độ} --- đặc trưng nào lệch xa mức bình thường thì token của nó mạnh lên. Nhờ vậy mô hình
biết con số đang đọc là bạch cầu hay huyết áp, thay vì nhận một dãy số vô danh.

Mỗi lần chụp vì thế cho một chuỗi 259 token: 257 token cột, một token tóm tắt cả hàng, và
một token vi sinh.

Vi sinh phải tách ra vì số chủng nuôi cấy được không cố định --- nhiều ca không mọc chủng
nào, một số ca mọc vài chủng, trung bình 1.7 chủng mỗi lần chụp. Đệm mọi ca về độ dài của
ca nhiều nhất thì gần như toàn bộ ô là rỗng, nên thay vào đó các chủng của một ca được
cộng gộp lại thành một token duy nhất, mỗi chủng có trọng số theo số lần nó mọc. Ba mươi
cột vi sinh dạng phẳng đã bị loại khỏi bảng để không đếm hai lần.


\subsubsection*{[D] \code{lớp1} --- điều kiện hoá bằng EHR ở mức patch}

Đây là đóng góp kiến trúc trung tâm của dự án. Patch làm \textbf{query}, token EHR làm
key/value --- tức là EHR \emph{điều biến biểu diễn ảnh} trước khi đầu đọc nhìn vào,
chứ không phải nối vào sau.

\begin{itemize}
\item Phép chiếu \code{pr} khởi tạo bằng \textbf{không} $\Rightarrow$ tại bước 0 mô hình
      đúng bằng bản ảnh thuần. Có assert kiểm điều này trước mỗi lượt chạy.
\item \code{ln\_a} chuẩn hoá \emph{đầu vào} của phép chiếu. Thiếu nó, arm đối chứng ở
      \code{38d} sụp từ $0.2416$ xuống $0.0681$.
\end{itemize}

\subsubsection*{[E] \code{lớp1Img} --- phim chụp trước ở mức patch}

Phim trước của cùng bệnh nhân được đưa vào ở mức patch, chạy trong không gian
\code{D\_CXR} trước phép chiếu \code{pin}. \textbf{Không} trừ thẳng
$\text{patch} - \text{prior}$ vì hai phim không được đăng ký không gian (tư thế khác,
độ hít vào khác).
Zero-init \emph{không} đủ để chặn một nhánh phụ lớn lên quá
nhanh. AdamW cập nhật khoảng $\text{lr}$ mỗi bước cho mỗi tham số bất kể gradient nhỏ
thế nào, nên nhánh xuất phát từ 0 vẫn lớn \emph{tuyến tính} theo bước. Đã sụp ba lần,
luôn ở epoch 2--4 tức ngay sau đỉnh warmup. Bản hiện tại dùng hệ số \textbf{LayerScale}
$10^{-4}$ có gradient thật, buộc mô hình phải \emph{học} để lớn lên. Ảnh nào không có
phim trước thì nhánh bị mask đi hẳn (28.7\% số ảnh).

\subsubsection*{[F] Hàm mất mát --- ASL}

Asymmetric Loss với $\gamma_{\text{neg}} = 4$, $\gamma_{\text{pos}} = 0$, clip $0.05$.

\subsubsection*{[G] Thành viên thứ hai --- mạng tích chập}

Thành viên thứ hai đổi backbone sang một mạng tích chập cỡ nhỏ, khởi tạo từ trọng số đã
huấn luyện trước trên \textbf{PadChest} thay vì từ ImageNet. Phần đầu đọc giữ nguyên như
pipeline chính: vẫn phim chụp trước ở mức patch, vẫn điều kiện hoá bằng EHR, vẫn bốn mươi
truy vấn lớp. Chỉ backbone là khác, và đó là chủ ý.

Trước mỗi lượt chạy, trọng số khởi tạo được đối chiếu với bản ImageNet gốc. Nếu hai bên
trùng khớp thì tệp PadChest đã không được nạp, và lượt chạy dừng lại thay vì âm thầm huấn
luyện từ khởi tạo sai.

CNN yếu hơn pipeline với RAD-DINO trước đó ($0.2463$ so với $0.2663$ trên val). Tuy nhiên
hai kiến trúc khác nhau nên trong quá trình huấn luyện chúng cập nhật theo hàm mất mát
theo những hướng khác nhau, và khi trộn lại thì sự đa dạng đó bổ sung cho nhau. Đó là lý do một thành viên yếu hơn vẫn làm cả ensemble mạnh lên.

\subsubsection*{[H] Trộn theo lớp}

Ensemble 20 file \code{t1\_prior\_both} $\oplus$ 16 file \code{cnn\_full}, trọng số
theo lớp \code{head 0.40 | medium 0.40 | tail 0.00}, cổng $0.85$. Trọng số được khoá
từ trước và đã kiểm lại sau R2a: tỉ lệ $r = \text{AP}_{\text{cnn}}/\text{AP}_{\text{ống}}$
gần như không đổi (head $0.981 \to 0.983$, medium $0.921 \to 0.923$, tail
$0.640 \to 0.653$), nên trọng số cũ vẫn đúng bối cảnh.

\paragraph{Tail nhận trọng số 0.} Tạm thời giữ nguyên như vậy vì CNN quá yếu ở tail để
đóng góp. Nghĩa là \textbf{16 nhãn tail hiện không nhận gì từ ensemble}, và đó là một
trong những dư địa rõ nhất còn lại.

\subsection{Hai arm huấn luyện}

Mỗi lượt huấn luyện chạy song song hai arm dùng \emph{chung một forward backbone},
nên arm đối chứng gần như miễn phí:

\begin{table}[h]
\centering
\small
\begin{tabular}{@{}llll@{}}
\toprule
\textbf{Arm} & \textbf{EHR} & \textbf{Phim trước} & \textbf{Vai trò} \\
\midrule
\code{mldec} & không & không & đối chứng ảnh thuần \\
\code{tpb}   & có    & có    & pipeline chính (\emph{t1\_prior\_both}) \\
\bottomrule
\end{tabular}
\caption{Hiệu số $G_{\text{ctx}} = \code{tpb} - \code{mldec}$ đo trực tiếp đóng góp của
ngữ cảnh EHR + phim trước, trong cùng một lượt chạy nên không có nhiễu giữa các lượt.}
\end{table}

\paragraph{Siêu tham số chốt.} $518^2$ $\cdot$ batch $32\times8 = 256$ hiệu dụng $\cdot$
10 epoch $\cdot$ LR $10^{-3}$ $\cdot$ WD $10^{-2}$ $\cdot$ warmup $5\%$ $\cdot$ clip $5.0$
$\cdot$ seed 42, 43. Thành viên CNN: batch $48\times6 = 288$, 8 epoch, LR $10^{-4}$.

% =====================================================================
\section{Các hướng đã đề ra: kết quả và trạng thái}
% =====================================================================

\subsection{Hướng đã hoàn thành và có lợi}

\begin{table}[h]
\centering
\small
\begin{tabular}{@{}l p{4.0cm}ll p{4.6cm}@{}}
\toprule
& \textbf{Hướng} & \textbf{Kết quả} & \textbf{Ghi chú} \\
\midrule
\XONG &
\textbf{R2a} --- gộp lát chọn vào train (74\% $\to$ 87\% dữ liệu có nhãn)
 & \textbf{+0.0042} test & 8/8 phép so lớp dương; CI95 $[+0.0033, +0.0051]$, loại 0 \\
\addlinespace
\XONG & Điều kiện hoá EHR ở mức patch (\code{lớp1})
 & $G_{\text{ctx}} \approx +0.0115$ & đóng góp ngữ cảnh, đo trong cùng lượt \\
\addlinespace
\XONG & Phim trước ở mức patch (\code{lớp1Img})
 & dương, nhỏ & LayerScale bắt buộc, nếu không sẽ sụp \\
\addlinespace
\XONG & Ensemble 2 kiến trúc, trộn theo lớp
 & \textbf{+0.0042} & cấu hình triển khai hiện tại \\
\bottomrule
\end{tabular}
\end{table}

\paragraph{Chi tiết R2a.} Lát chọn (39{,}086 ảnh) vốn là thước đo để chọn siêu tham số.
Khi mọi lựa chọn đã khoá, nó hết vai trò và được thu hồi vào train.

\begin{table}[h]
\centering
\small
\begin{tabular}{@{}lrrrrr@{}}
\toprule
\textbf{Cấu hình} & \textbf{n file} & \textbf{macro} & \textbf{head} & \textbf{medium} & \textbf{tail} \\
\midrule
\code{53/r3} pipeline cũ   & 20 & 0.2641 & 0.4758 & 0.1942 & 0.1155 \\
\code{64/r2a} pipeline R2a & 20 & \textbf{0.2663} & 0.4775 & 0.1958 & \textbf{0.1193} \\
\code{41e/r1} cnn cũ        & 16 & 0.2435 & 0.4667 & 0.1788 & 0.0739 \\
\code{64b/r2a} cnn R2a      & 16 & \textbf{0.2463} & 0.4695 & 0.1808 & \textbf{0.0779} \\
\bottomrule
\end{tabular}
\caption{Val, thang lớp \emph{cũ} (13/16/11). pipeline $+0.0022$ macro và $+0.0039$ tail;
CNN $+0.0028$ macro và $+0.0040$ tail. Cơ chế: thêm $17.8\%$ dữ liệu là thêm $17.8\%$ ca
dương, mà nhãn hiếm nhất chỉ có khoảng 150 ca dương nên thêm vài chục ca là đáng kể;
head thì đã bão hoà.}
\end{table}

\subsection{Hướng đã ĐÓNG}
\label{sec:closed}

Mỗi dòng dưới đây là một giả thuyết đã bị bác bỏ bằng phép đo, không phải bị bỏ vì hết
thời gian. Ghi lại đầy đủ vì kết quả âm cũng là thông tin nghiên cứu.

\begin{longtable}{@{}l p{4.5cm} p{9.2cm}@{}}
\toprule
& \textbf{Hướng} & \textbf{Bằng chứng đóng} \\
\midrule
\endhead
\DONG &Đưa EHR vào sâu trong backbone &
Đo chênh lệch giữa nhánh có điều kiện hoá sâu và nhánh đối chứng chạy song song: chênh
lệch \textbf{âm}, lặp lại ổn định qua nhiều seed. Đưa EHR vào sâu kém hơn đưa vào ở đầu đọc. \\
\addlinespace
\DONG &Mở khoá bốn khối trên cùng của backbone &
Điểm giảm $0.007$ đến $0.010$. Hai seed cho cùng kết luận, và nhánh đối chứng chạy song
song vẫn khoẻ, nên đây không phải sự cố của một lượt chạy. \\
\addlinespace
\DONG &Huấn luyện lại riêng tầng phân loại &
Ý tưởng là tách tầng cuối ra học lại trên dữ liệu đã cân bằng. Đo ra thì tầng này
\textbf{đã hội tụ}, không còn gì để tách. \\
\addlinespace
\DONG &Ghép đặc trưng đa tỉ lệ từ mạng tích chập &
Trượt cổng. Nguyên nhân đo được ở \code{62}: tầng thứ ba gần như chỉ mang \emph{một}
chiều thông tin thực (hạng hiệu dụng $1.1$ trên 320 kênh), nên 1024 khoá thêm vào gần
như trùng nhau và chỉ làm loãng chú ý. \\
\addlinespace
\DONG &Tăng độ phân giải chỉ ở lúc suy luận &
Ảnh vào lớn hơn nhưng đầu đọc vẫn là bản đã huấn luyện ở 518. Số lượng vùng ảnh thay đổi
khiến đầu vào lệch khỏi phân phối lúc huấn luyện: $-0.0028$. \\
\addlinespace
\DONG &Thêm seed trên cùng một backbone đông cứng &
$+0.0008$, dưới sàn nhiễu $0.0007$. Điều này \emph{không} đóng hướng huấn luyện lại
backbone độc lập rồi mới trộn --- đó là phép đo khác hẳn. \\
\addlinespace
\DONG &Hiện tượng co cụm biểu diễn lớp &
Kiểm trực tiếp giả thuyết biểu diễn của mỗi lớp co về một điểm. Không quan sát thấy. \\
\addlinespace
\DONG &Trọng số trộn phụ thuộc chất lượng từng nhãn &
\textbf{Tương quan giả do nhiễu dùng chung.} Đo hai đại lượng trên \emph{cùng} một tập
cho tương quan $+0.598$; đo trên hai tập \emph{tách rời} thì tụt còn $+0.136$ với
$p = 0.40$. Cả hai đại lượng đều chứa điểm của pipeline, nên nhiễu của nó tự sinh ra
tương quan. \\
\addlinespace
\DONG &Dùng độ tương đồng giữa hai thành viên làm biến quyết định &
Đại lượng này đo được rất ổn định, nhưng \textbf{không dự báo được gì}: tương quan của
nó với mức cải thiện thực tế chỉ $+0.017$. \\
\addlinespace
\DONG &Bốn cách trộn rẻ, không cần GPU &
Trung bình thứ hạng thay cho trung bình xác suất, quét lại trọng số theo lớp, cân theo
chất lượng từng epoch, và thêm thành viên thứ ba. \textbf{Tất cả cho $+0.0000$.} Cấu
hình đang dùng vẫn tốt nhất ($0.2709$); thêm thành viên thứ ba còn tụt xuống $0.2693$. \\
\addlinespace
\DONG &Cắt vùng nghi ngờ rồi phóng to &
Vùng cắt trung bình chiếm \textbf{79\% diện tích ảnh}, nên độ phóng đại thực chỉ
$1.06$ lần và kết quả $+0.0003$ không có ý nghĩa. Cắt hẹp hơn cũng không cứu được, vì
dấu hiệu bệnh trải khắp lồng ngực chứ không gom vào một chỗ. \\
\addlinespace
\DONG &Định vị vùng mô hình đang nhìn &
Che bớt ảnh, chỉ chừa lại số ít vùng mà mô hình chú ý nhiều nhất. Giữ $5\%$ số vùng vẫn
đạt $96$--$99\%$ điểm gốc; giữ $20\%$ thì nhóm nhãn hiếm còn \emph{cao hơn} cả lượt dùng
toàn ảnh. Mô hình không hề nhìn nhầm chỗ --- và điều đó lấy đi luôn lý do tồn tại của
hướng cắt vùng rồi phóng to. \\
\addlinespace
\DONG &Đưa cả pipeline lên độ phân giải 1022 &
Điểm chung giảm $0.0055$, nhóm nhãn hiếm giảm $0.0138$. Mọi nhãn thiệt hại nặng đều
thuộc nhóm hiếm, mọi nhãn hưởng lợi đều thuộc nhóm trung bình. \\
\addlinespace
\DONG &Gắn lại nhãn bằng mô hình ngôn ngữ &
Nhãn của tập test bẩn theo \emph{đúng cùng một kiểu} với nhãn huấn luyện, và lỗi đến từ
bộ phân tích văn bản nên mang tính hệ thống chứ không ngẫu nhiên. Dạy mô hình đoán đúng
bệnh thật sẽ kéo nó ra \emph{xa} thước đo đang được chấm. \\
\bottomrule
\end{longtable}


\paragraph{ViT thua ở độ phân giải cao.} Bốn mươi class query phải học phân bố
attention trên $5{,}329$ Key thay vì $1{,}369$ --- gấp $3.9$ lần. Nhãn có 26{,}000 ca
dương thì học được; nhãn có 129 ca dương thì \textbf{pha loãng thống kê}. Chỉ dấu phụ:
$G_{\text{ctx}}$ tụt từ $+0.0115$ xuống $+0.0029$.

\begin{table}[h]
\centering
\small
\begin{tabular}{@{}lrrr@{}}
\toprule
& \textbf{64 @518} & \textbf{68 @1022} & \textbf{chênh} \\
\midrule
macro  & 0.2645 & 0.2590 & $-0.0055$ \\
head   & 0.5593 & 0.5545 & $-0.0048$ \\
medium & 0.2227 & 0.2256 & $\mathbf{+0.0029}$ \\
tail   & 0.1379 & 0.1241 & $\mathbf{-0.0138}$ \\
\bottomrule
\end{tabular}
\caption{Nhãn hại nặng nhất: \code{Clavicle Fracture} $-0.0916$ (prevalence $0.065\%$).
Nhãn lợi nhất: \code{Granuloma} $+0.0484$ ($1.145\%$).}
\end{table}

Khi chạy thí nghiệm về độ phân giải thì mới chỉ chạy 3 epoch, 1 seed, $E=6$.
Đây là phép \emph{sàng}, không phải phép \emph{đóng} tuyệt đối. Phát biểu an toàn: \emph{ở
cùng mức chi phí, ViT không đổi được độ phân giải thành lợi ích}. Việc CNN có làm được hay
không là câu hỏi riêng, đang chờ \code{nb68b} trả lời.

\subsection{Hướng còn MỞ}

\begin{center}
\begin{tabular}{@{}l l r@{}}
\toprule
& \textbf{Hướng} \\
\midrule
\MO & \textbf{\code{nb68b}} --- ConvNeXtV2 @1022  \\
\MO & \textbf{Weighted asymmetric loss}  \\
\MO & K-fold + fine-tune backbone theo từng fold  \\
\MO & Sinh ảnh diffusion cho nhãn hiếm \\
\bottomrule
\end{tabular}
\end{center}


\end{document}